% ---
% id: transmitter/rem:segment-scope
% kind: rem
% title: "Scope of the source-segment conditions"
% summary: >-
%   The registry conditions decide that a link names a claim proved for its
%   kind, carries that claim's assumptions and comes from the admitting
%   module; they do not decide that the claim applies to these descriptors.
%   Composing the segment with the receiver segment is stated elsewhere.
% prominence: standard
% status: accepted
% uses: [def:source-segment]
% mentions: [def:source-link, reductions/lem:precomposition, reductions/thm:two-sided, reductions/def:chain, protocol-compiler/def:chaincond]
% assumes: []
% refs: []
% provenance: newly-derived
% lean:
% ---

\begin{remark}\label{rem:segment-scope}
\emph{What (A1) to (A3) decide.} The conditions (A1) to (A3) of
Definition~\ref{def:source-segment} decide that a link's obligation names a
claim of this repository that was proved for its kind, that the assumptions
that obligation carries --- the claim's closure and the attestation for its
model --- travel in the ledger, and that the segment came from the module that
admits it. They do not decide that the named claim's conclusion applies to
these descriptors: family membership is asserted by the emitting module and
recorded as the fourth component of the obligation record
(Definition~\ref{def:source-link}\,(c2), (d)), and no consumer re-proves a
theorem. A link's obligation thus has two halves, what a consumer decides and
what it takes from the emitting module. The second is exactly what rests on
the emitting module for every artefact this repository passes between
modules, and it is the smaller half: the registry closes the case of a link
naming a claim that proves nothing about its kind, which is the case no other
condition reaches.

\emph{What is decided elsewhere.} The definition admits a \emph{kind} as a
type, on the strength of a claim and a registry entry, and it decides of a
presented segment the conditions (S1) to (S4) and (A1) to (A3). It asserts of
no specification that it declares a non-empty segment: whether one does is
\emph{data} of the presented chain, read off the segment itself. Two things a
non-empty segment needs are not the definition's and are named rather than
restated: the rule by which the segment composes with a receiver segment,
which is \xref{reductions/lem:precomposition} and
\xref{reductions/thm:two-sided} with the admission rule of
\xref{reductions/def:chain}, in \tref{trusted-source mode}; and the consumer that
decides that rule on an assembled declaration, which is
\xref{protocol-compiler/def:chaincond}. That consumer also refuses a
non-empty segment in \tref{untrusted-source mode}.
\end{remark}
